\documentclass[11pt]{article}

\usepackage[a4paper,margin=1in]{geometry}
\usepackage{amsmath,amssymb,amsthm}
\usepackage{mathtools}
\usepackage{hyperref}
\usepackage{xcolor}
\usepackage{listings}
\usepackage{booktabs}
%\usepackage{microtype}
\usepackage[T1]{fontenc}
\usepackage{bussproofs}
\EnableBpAbbreviations

\let\AXM\AXC
\let\BIM\BIC
\let\UIM\UIC
\let\TIM\TIC
\let\RLM\RL

% --- Theorem environments ---
\newtheorem{theorem}{Theorem}
\newtheorem{lemma}{Lemma}
\newtheorem{definition}{Definition}
\newtheorem{remark}{Remark}

% --- Listings setup for Lean ---
\lstdefinelanguage{Lean}{
  morekeywords={
    import, namespace, end, inductive, def, theorem, abbrev,
    by, intro, apply, exact, simp, constructor, cases, with,
    match, open, have, let
  },
  sensitive=true,
  morecomment=[l]{--},
  morestring=[b]"
}

% Unicode rendering to math for readability (non-Lean readers)
\lstset{
  language=Lean,
  basicstyle=\ttfamily\small,
  keywordstyle=\color{blue!60!black},
  commentstyle=\color{green!40!black},
  stringstyle=\color{red!50!black},
  showstringspaces=false,
  breaklines=true,
  frame=single,
  columns=fullflexible,
  mathescape=true,
  literate=
    {⟶}{{$\to$}}1
    {~}{{$\neg$}}1
    {∅}{{$\emptyset$}}1
    {⊢}{{$\vdash$}}1
    {⊧}{{$\models$}}1
    {⋎}{{$\lor$}}1
    {∼}{{$\sim$}}1
    {¬}{{$\neg$}}1
    {∀}{{$\forall$}}1
    {→}{{$\to$}}1
    {↔}{{$\leftrightarrow$}}1
    {Γ}{{$\Gamma$}}1
    {φ}{{$\varphi$}}1
    {ψ}{{$\psi$}}1
    {∈}{{$\in$}}1
    {⊥}{{$\bot$}}1
    {∧}{{$\land$}}1
    {∨}{{$\lor$}}1
}


% --- Notation in the paper ---
\newcommand{\Imp}{\rightarrow}
\newcommand{\Not}{\neg}
\newcommand{\Or}{\vee}
\newcommand{\Sat}{\models}
\newcommand{\Prov}{\vdash}
\newcommand{\llbracket}{[\![}
\newcommand{\rrbracket}{]\!]}

\title{Completeness of Propositional Logic in Lean 4:\\ A Verified Bridge from Semantics to a Natural-Deduction Interface}
\author{\textbf{Author Name} \\ Undergraduate Thesis \\ University \\ \texttt{email@example.com}}
\date{January 2026}

\begin{document}
\maketitle



\begin{abstract}
This chapter develops a machine-checked verification pipeline for classical propositional logic in Lean~4.
The construction is deliberately modular: we define a surface language with only atoms, negation, and implication;
give its standard classical semantics; and formalize a verified translation into the \emph{Foundation} library's propositional calculus.
A bridge lemma (semantics preservation) is proved, allowing us to import the library's completeness result.
Finally, we package the result as a \emph{Provability Interface}: a genuine Natural Deduction system with contexts and discharge,
augmented by a verified oracle that lifts library proofs to the surface.
The exposition targets readers new to Lean, mapping code to mathematics line-by-line, and explicitly outlining
the architectural separation between the surface interface and the library backend.
\end{abstract}

\tableofcontents

\section{Motivation and relation to Stansifer (2001)}
\label{sec:stansifer}

The guiding inspiration for this chapter is Ryan Stansifer's paper
\emph{Completeness of Propositional Logic as a Program}~\cite{Stansifer2001}.
Stansifer emphasizes a viewpoint that has become central in modern mechanized reasoning:
a constructive completeness proof can be executed as a program that takes a tautology $\varphi$ and returns a
derivation of $\varphi$ in a fixed proof system.

In Stansifer's development, propositional formulas use only $\Not$ and $\Imp$,
and proofs are explicit proof trees in a Hilbert system (specifically \L{}ukasiewicz axioms plus modus ponens).
The program produces proof objects and is correct by construction.
The same conceptual foundation also appears in proof theory and type theory:
\emph{proof objects are computational artifacts} and can be manipulated systematically.

Lean provides a setting where this slogan is made precise:
proofs are typed terms, and the verification of correctness is delegated to a small trusted kernel.
This chapter adopts Stansifer's conceptual perspective while exploiting Lean's ecosystem:
instead of generating large explicit proof trees, the key engineering move is a verified translation into a library
calculus that already contains a completeness theorem.

\section{Why Lean is a good fit (a short primer)}
\label{sec:lean-primer}

This section explains what Lean is checking, why the approach is trustworthy, and how to read the code as mathematics.

% ============================================================
% PREPARATORY SECTION (insert after the Lean primer section)
% This section is adapted in spirit from the documentation excerpt
% you provided, but rewritten in a single, thesis-consistent voice,
% and connected directly to *your* Lean development.
%
% IMPORTANT PREAMBLE ADDITIONS (add near your other \usepackage lines):
%   \usepackage{bussproofs}
%   \EnableBpAbbreviations
%
% bussproofs provides \AXM, \BIM, \UIM, \TIM, \RLM, \DP used for proof trees.
% ============================================================

\section{Preparatory background: natural deduction rules and their Lean realizations}
\label{sec:prep-doc}

This section prepares the reader to interpret Lean proofs as natural deduction derivations.
The emphasis is on the \emph{structural correspondence} between the familiar proof trees of
symbolic logic and the Lean expressions that inhabit the corresponding proposition types.

\subsection{From logic puzzles to formal rules (why proof trees)}
A standard logic puzzle (such as ``Malice and Alice'') is solved by organizing informal reasoning into
reusable patterns: ``if $\dots$ then $\dots$,'' ``either $\dots$ or $\dots$,'' ``assume $\dots$ and derive a contradiction,'' and so on.
Symbolic logic isolates these patterns by replacing content-rich statements with propositional variables
$A,B,C,\dots$ and replacing English connectives by formal connectives $\wedge,\vee,\neg,\to$.
The crucial payoff is that \emph{inference} becomes the study of a small number of \emph{rules}.

We will represent rules as proof trees in the standard way:
if $C$ is derivable from $A$ and $B$, we draw
\[
\frac{A \qquad B}{C}.
\]
In natural deduction, these trees can contain \emph{subproofs} with temporarily assumed hypotheses
that are later \emph{discharged} (canceled). This is the formal counterpart of hypothetical reasoning in prose:
\begin{quote}
``Assume $A$. Then show $B$. Therefore $A\to B$.'' 
\end{quote}

\subsection{The central idea: propositions as types, proofs as terms}
Lean makes the proof-tree discipline both precise and executable:
\[
\text{a proposition }P \text{ is a type}, \qquad \text{a proof of }P \text{ is a term of type }P.
\]
Thus, in Lean, proving $P$ means constructing a term whose type is $P$.
This is not metaphorical: the Lean kernel checks the term is well-typed,
and that check is the formal correctness criterion for the proof.

In particular, implication is realized literally as a function type:
\[
A \to B \quad\text{is interpreted as the type of functions }(A \to B).
\]
This is why implication elimination (modus ponens) becomes ordinary function application,
and implication introduction becomes lambda abstraction (\texttt{fun} or $\lambda$).

\subsection{Rules of inference as proof trees}
We now list the most important rules (in the proof-tree style used in the documentation)
and immediately show the corresponding Lean patterns. The point is not to learn syntax by memorization,
but to see that the Lean expression \emph{is the proof tree}.

\subsubsection{Implication elimination (modus ponens)}
The elimination rule for implication is:

\begin{prooftree}
\AXM{$A \to B$}
\AXM{$A$}
\RLM{$\to\mathrm{E}$}
\BIM{$B$}
\end{prooftree}

In Lean, if \texttt{h1 : A $\to$ B} and \texttt{h2 : A}, then \texttt{h1 h2 : B}.
That is, modus ponens is function application.

\begin{lstlisting}
section
  variable (A B : Prop)
  variable (h1 : A → B) (h2 : A)

  #check h1 h2  -- has type B
end
\end{lstlisting}

\subsubsection{Implication introduction (hypothetical reasoning)}
The introduction rule is:
\begin{center}
\begin{prooftree}
\AXM{}
\RLM{1}
\UIM{$A$}
\noLine
\UIM{$\vdots$}
\noLine
\UIM{$B$}
\RLM{$1\;\;\to\mathrm{I}$}
\UIM{$A \to B$}
\end{prooftree}
\end{center}
The label ``1'' indicates the hypothesis $A$ is discharged at the end. In Lean, this is exactly a lambda term:
if assuming \texttt{h : A} we can build a term of type \texttt{B}, then \texttt{fun h : A $\mapsto$ ...} has type \texttt{A $\to$ B}.

\begin{lstlisting}
section
  variable (A B : Prop)
  -- a skeleton: to prove A → B, we assume A and prove B
  example : A → B :=
  fun (h : A) =>
    show B from by
      -- here we can use h : A to derive B
      sorry
end
\end{lstlisting}

\subsubsection{Conjunction}
The introduction rule:
\begin{center}
\begin{prooftree}
\AXM{$A$}
\AXM{$B$}
\RLM{$\wedge\mathrm{I}$}
\BIM{$A \wedge B$}
\end{prooftree}
\end{center}
is \texttt{And.intro} in Lean. The elimination rules:
\begin{center}
\begin{prooftree}
\AXM{$A \wedge B$}
\RLM{$\wedge\mathrm{E_l}$}
\UIM{$A$}
\end{prooftree}
\qquad
\begin{prooftree}
\AXM{$A \wedge B$}
\RLM{$\wedge\mathrm{E_r}$}
\UIM{$B$}
\end{prooftree}
\end{center}
are \texttt{And.left} and \texttt{And.right}.

\begin{lstlisting}
section
  variable (A B : Prop)
  variable (hA : A) (hB : B)

  #check And.intro hA hB  -- : A ∧ B

  variable (h : A ∧ B)
  #check And.left h   -- : A
  #check And.right h  -- : B
end
\end{lstlisting}

\subsubsection{Disjunction and proof by cases}
Disjunction introduction:
\begin{center}
\begin{prooftree}
\AXM{$A$}
\RLM{$\vee\mathrm{I_l}$}
\UIM{$A \vee B$}
\end{prooftree}
\qquad
\begin{prooftree}
\AXM{$B$}
\RLM{$\vee\mathrm{I_r}$}
\UIM{$A \vee B$}
\end{prooftree}
\end{center}
corresponds to \texttt{Or.inl} and \texttt{Or.inr}. Disjunction elimination is the formal case split:
\begin{center}
\begin{prooftree}
\AXM{$A \vee B$}
\AXM{}
\RLM{1}
\UIM{$A$}
\noLine
\UIM{$\vdots$}
\noLine
\UIM{$C$}
\AXM{}
\RLM{2}
\UIM{$B$}
\noLine
\UIM{$\vdots$}
\noLine
\UIM{$C$}
\RLM{$\vee\mathrm{E}$}
\TIM{$C$}
\end{prooftree}
\end{center}
In Lean, \texttt{Or.elim} expresses exactly this two-branch structure:
provide a proof of $C$ assuming $A$, and a proof of $C$ assuming $B$.

\begin{lstlisting}
section
  variable (A B C : Prop)
  variable (h : A ∨ B)
  variable (ha : A → C)
  variable (hb : B → C)

  example : C :=
    Or.elim h
      (fun a : A => ha a)
      (fun b : B => hb b)
end
\end{lstlisting}

\subsubsection{Negation, falsity, and \emph{ex falso}}
Negation is treated in Lean as implication into \texttt{False}:
\[
\neg A \quad \equiv\quad A \to \bot.
\]
Thus the negation elimination rule
\begin{center}
\begin{prooftree}
\AXM{$\neg A$}
\AXM{$A$}
\RLM{$\neg\mathrm{E}$}
\BIM{$\bot$}
\end{prooftree}
\end{center}
is again function application: if \texttt{hnot : $\neg$A} and \texttt{hA : A} then \texttt{hnot hA : False}.
And from falsity we may derive anything (\emph{ex falso}):
\begin{center}
\begin{prooftree}
\AXM{$\bot$}
\RLM{$\bot\mathrm{E}$}
\UIM{$A$}
\end{prooftree}
\end{center}
implemented as \texttt{False.elim}.

\begin{lstlisting}
section
  variable (A : Prop)
  variable (hFalse : False)

  -- ex falso: from False, derive anything
  #check False.elim hFalse  -- : A

  variable (hnotA : ¬ A) (hA : A)
  #check hnotA hA  -- : False
end
\end{lstlisting}

\subsection{How this preparatory material connects to the chapter}
The rest of the chapter specializes these general principles to the propositional fragment we use in the implementation.

\subsubsection{Why the surface language uses only \texorpdfstring{$\neg$ and $\to$}{negation and implication}}
Stansifer's paper uses only $\neg$ and $\to$ as primitives, treating other connectives as definable.
We do the same, for two reasons:
\begin{enumerate}
\item It keeps the surface syntax small and the semantics easy to explain.
\item It makes the bridge to the Foundation library conceptually clean:
      implication is translated using the classical identity
      \[
      (A \to B) \leftrightarrow (\neg A \vee B).
      \]
\end{enumerate}

\subsubsection{Where the proof trees appear in our Lean development}
In our files, proof trees appear in two different but tightly related ways:

\paragraph{(1) Proof trees as inductive objects.}
When we define an inductive predicate like
\texttt{ND : Formula $\to$ Prop}, each constructor is an inference rule.
A term \texttt{p : ND $\varphi$} is a structured proof object.
This is the same viewpoint as Stansifer's proof trees, except Lean checks correctness by typing, not by runtime checks.

\paragraph{(2) Proof trees as terms built by tactics.}
When we write
\texttt{theorem eval\_tr ... := by ...},
Lean’s tactic mode constructs a term of the required type.
The tactics can be understood as a readable script that produces the proof tree.

\subsection{A final ``bridge'' example: why implication corresponds to \texorpdfstring{$\neg A \vee B$}{neg A or B}}
The translation step of this thesis uses the classical equivalence
\[
(A \to B) \leftrightarrow (\neg A \vee B).
\]
This identity is exactly what makes the surface semantics align with the target semantics.
At the level of proof trees, this is the reason that the implication case of our bridge lemma reduces to
a case split on whether $A$ holds:
\[
A \to B \quad\text{is true iff}\quad (\neg A)\ \text{or}\ B.
\]
In Lean, the corresponding proof is the combination of:
\begin{itemize}
\item a case split (\texttt{by\_cases} / \texttt{cases}) on $A$,
\item and a disjunction introduction (\texttt{Or.inl} / \texttt{Or.inr}) to build $\neg A \vee B$.
\end{itemize}

This is precisely the same logical structure as the informal puzzle reasoning in the documentation:
\emph{split into exhaustive cases, derive the same conclusion, and conclude the goal holds.}

\subsection{Takeaway}
Lean does not replace mathematics with syntax; it enforces that the mathematics has been fully specified.
Proof trees, hypotheses, discharge, and case splits all appear explicitly:
\begin{itemize}
\item as inductive constructors (when we encode a proof system), and
\item as proof terms (when we build proofs directly).
\end{itemize}
With this perspective in place, the remaining chapter can be read as a sequence of proof-tree constructions
whose correctness is checked automatically by the Lean kernel.

\subsection{Proofs as typed terms (Curry--Howard)}
Lean is based on dependent type theory. The key operational principle is the Curry--Howard correspondence:
\[
\text{Propositions} \ \leftrightarrow\  \text{Types},
\qquad
\text{Proofs} \ \leftrightarrow\  \text{Terms (programs)}.
\]
Under Curry--Howard, an implication $P\Imp Q$ corresponds to a function type $P \to Q$:
a proof of $P\Imp Q$ is a function that transforms proofs of $P$ into proofs of $Q$.
This relationship is the conceptual bridge between natural deduction and typed functional programming.
For background and a systematic treatment, see Sørensen--Urzyczyn~\cite{SorensenUrzyczyn2006} and Pierce~\cite{PierceTAPL2002}.%

\subsection{Inductive definitions as inference systems}
Lean's \texttt{inductive} definitions encode the \emph{least} set closed under specified constructors.
This is exactly how inference rules generate provability relations.
Thus, an inductive predicate
\[
\mathrm{ND}(\varphi) : \mathsf{Prop}
\]
with constructors corresponding to axiom schemata and modus ponens gives a fully formal notion of derivability.
Any Lean term of type \texttt{ND $\varphi$} is, by construction, a derivation of $\varphi$ according to those rules.

\subsection{Computation inside proofs}
Lean definitions are computational: they can be unfolded and reduced.
The \texttt{simp} tactic performs definitional unfolding and rewriting using registered lemmas.
In this chapter, \texttt{simp} is used primarily to expand the recursive definitions of semantics and translation,
reducing a proof to elementary propositional reasoning that is then discharged explicitly.

\subsection{Libraries as verified mathematical infrastructure}
A distinctive advantage of Lean is that large theorems can be reused safely:
the system checks that a theorem is applied with correct hypotheses and correct types.
In this chapter, the Foundation library provides a completeness theorem for a propositional calculus with respect to classical semantics;
we supply the interface (translation + semantics preservation) that makes the theorem applicable to our surface language.

\section{Glossary (Lean term \texorpdfstring{$\leftrightarrow$}{<->} mathematical meaning)}
\label{sec:glossary}

\begin{center}
\begin{tabular}{@{}ll@{}}
\toprule
\textbf{Lean notation / identifier} & \textbf{Mathematical meaning} \\
\midrule
\texttt{Formula} & Surface formula syntax $\varphi ::= p \mid \Not\varphi \mid (\varphi\Imp\psi)$ \\
\texttt{Valuation := String $\to$ Prop} & Classical valuation $v : \mathsf{Var} \to \{\mathsf{true},\mathsf{false}\}$ (encoded as $\mathsf{Prop}$) \\
\texttt{eval v $\varphi$} & Semantics $\llbracket \varphi \rrbracket_v$ \\
\texttt{IsTautology $\varphi$} & $\forall v,\ \llbracket \varphi \rrbracket_v$ \\
\texttt{F := NNFormula String} & Target formula type in Foundation \\
\texttt{tr : Formula $\to$ F} & Translation $t(\varphi)$ into Foundation syntax \\
\texttt{v $\models$ $\psi$} / \texttt{Sat v $\psi$} & Library satisfaction relation $v \Sat \psi$ \\
\texttt{T $\vdash$ $\psi$} & Library provability $T \Prov \psi$ \\
\texttt{eval\_tr} & Bridge lemma: $\llbracket\varphi\rrbracket_v \leftrightarrow (v\Sat t(\varphi))$ \\
\texttt{provable\_tr\_of\_tautology} & Tautology $\Rightarrow$ library provability of translation \\
\texttt{ND $\Gamma$ $\varphi$} & Genuine ND proof of $\varphi$ from assumptions $\Gamma$ \\
\texttt{simp} & Definitional unfolding + rewriting (normalization by simplification) \\
\texttt{induction $\varphi$} & Structural induction on the syntax of $\varphi$ \\
\bottomrule
\end{tabular}
\end{center}

\section{Surface syntax and semantics (\texttt{Syntax.lean})}
\label{sec:syntax}

\subsection{Mathematics}
We fix a set of variables $\mathsf{Var}$ and define formulas by:
\[
\varphi ::= p \mid \Not\varphi \mid (\varphi\Imp\psi).
\]
A valuation $v$ assigns truth-values to atoms. We define evaluation:
\[
\llbracket p \rrbracket_v = v(p),\quad
\llbracket \Not\varphi \rrbracket_v = \Not \llbracket \varphi \rrbracket_v,\quad
\llbracket \varphi\Imp\psi \rrbracket_v = (\llbracket\varphi\rrbracket_v \Imp \llbracket\psi\rrbracket_v).
\]
A tautology is a formula true under all valuations:
\[
\mathrm{Taut}(\varphi) \iff \forall v,\ \llbracket \varphi \rrbracket_v.
\]

\subsection{Proof sketch (file-to-math map)}
\begin{itemize}
\item \texttt{inductive Formula} implements the grammar above.
\item \texttt{def eval} implements $\llbracket\cdot\rrbracket_v$ by recursion on syntax.
\item \texttt{def IsTautology} implements $\forall v,\ \llbracket\varphi\rrbracket_v$.
\end{itemize}

\subsection{Lean code (excerpt)}
\begin{lstlisting}
namespace Thesis.Prop

inductive Formula : Type
| atom : String → Formula
| neg  : Formula → Formula
| impl : Formula → Formula → Formula
deriving DecidableEq, Repr

prefix:60 "~" => Formula.neg
infixr:55 " ⟶ " => Formula.impl

abbrev Valuation : Type := String → Prop

def eval (v : Valuation) : Formula → Prop
| .atom s   => v s
| .neg p    => ¬ eval v p
| .impl p q => eval v p → eval v q

def IsTautology (φ : Formula) : Prop :=
  ∀ v : Valuation, eval v φ

end Thesis.Prop
\end{lstlisting}

\section{Translation into Foundation (\texttt{BridgeToFoundation.lean})}
\label{sec:translation}

\subsection{Mathematics}
We translate implication into disjunction and negation using the classical equivalence:
\[
(\varphi \Imp \psi) \equiv (\Not\varphi \Or \psi).
\]
Thus we define:
\[
t(p)=p,\qquad t(\Not\varphi)=\Not t(\varphi),\qquad t(\varphi\Imp\psi)=\Not t(\varphi)\Or t(\psi).
\]
This choice is standard in proof theory and matches the truth-table semantics for implication.

\subsection{Proof sketch (file-to-math map)}
\begin{itemize}
\item \texttt{abbrev F} fixes the target formula type (Foundation's \texttt{NNFormula}).
\item \texttt{def tr} implements $t(\cdot)$ with the key implication clause $t(\varphi\Imp\psi)=\Not t(\varphi)\Or t(\psi)$.
\item \texttt{abbrev Sat} names the library satisfaction predicate $v\Sat\Phi$.
\end{itemize}

\subsection{Lean code (excerpt)}
\begin{lstlisting}
import Foundation.Propositional.ClassicalSemantics.Tait
import Thesis.Prop.Syntax

open Classical

namespace Thesis.Prop

abbrev F : Type := LO.Propositional.NNFormula String

def tr : Formula → F
| .atom s   => LO.Propositional.NNFormula.atom s
| .neg p    => (∼ (tr p))
| .impl p q => (∼ (tr p)) ⋎ (tr q)

abbrev Sat (v : Valuation) (ψ : F) : Prop :=
  v ⊧ ψ

end Thesis.Prop
\end{lstlisting}

\section{The bridge lemma (\texorpdfstring{\texttt{eval\_tr}}{eval\_tr}) and why it is the key}
\label{sec:bridge}

\subsection{Statement}
\begin{lemma}[Semantics preservation]\label{lem:bridge}
For every valuation $v$ and surface formula $\varphi$,
\[
\llbracket \varphi \rrbracket_v \;\Longleftrightarrow\; v \Sat t(\varphi).
\]
\end{lemma}

\subsection{Proof idea}
The proof is by structural induction on $\varphi$.
The only substantive step is the implication case: after unfolding definitions,
the goal reduces to the classical tautology
\[
(P\Imp Q) \leftrightarrow (\Not P \Or Q),
\]
which is proved by a standard case split on $P$.

\subsection{Proof sketch (line-by-line structure)}
\begin{itemize}
\item \texttt{induction $\varphi$} corresponds to induction on the grammar of formulas.
\item In the atom case, \texttt{simp} reduces both sides to the same proposition $v(p)$.
\item In the negation case, \texttt{simp} reduces the goal to the induction hypothesis under $\Not$.
\item In the implication case, \texttt{simp} reduces to a propositional goal and the proof proceeds by:
  \begin{enumerate}
    \item assuming $P\Imp Q$ and proving $\Not P\Or Q$ via cases on $P$;
    \item assuming $\Not P\Or Q$ and proving $P\Imp Q$ by assuming $P$ and eliminating the disjunction.
  \end{enumerate}
\end{itemize}

\subsection{Lean code (as implemented)}
\begin{lstlisting}
theorem eval_tr (v : Valuation) : ∀ φ : Formula, eval v φ ↔ Sat v (tr φ) := by
  intro φ
  induction φ with
  | atom s => simp [eval, tr, Sat]
  | neg p ih =>
      simp [eval, tr, Sat, ih]
  | impl p q ihp ihq =>
      simp [eval, tr, Sat, ihp, ihq]
      tauto
\end{lstlisting}

\section{Completeness via Foundation (\texttt{CompletenessViaFoundation.lean})}
\label{sec:completeness}

\subsection{Mathematics}
The Foundation library provides a completeness theorem of the form:
\[
T \models \Phi \Longrightarrow T \Prov \Phi,
\]
i.e.\ semantic consequence implies provability (for the library calculus).
Specializing to the empty theory $T=\emptyset$:
\[
(\forall v,\ v\Sat \Phi)\Longrightarrow \emptyset \Prov \Phi.
\]
Now if $\varphi$ is a tautology in the surface semantics, then by Lemma~\ref{lem:bridge} we have
$\forall v,\ v\Sat t(\varphi)$, hence $\emptyset\Prov t(\varphi)$.

\subsection{Proof sketch (file-to-math map)}
\begin{itemize}
\item \texttt{provable\_tr\_of\_tautology} is the composition:
  \[
  \mathrm{Taut}(\varphi) \xRightarrow{\text{bridge lemma}} \forall v,\ v\Sat t(\varphi)
  \xRightarrow{\text{library completeness}} \emptyset\Prov t(\varphi).
  \]
\item In the Lean proof:
  \begin{enumerate}
  \item \texttt{apply ... completeness!} invokes the library theorem;
  \item the remaining goal is to show semantic consequence from tautologyhood;
  \item \texttt{exact (eval\_tr v $\varphi$).1 ...} converts surface truth into library satisfaction.
  \end{enumerate}
\end{itemize}

\subsection{Lean code (as implemented)}
\begin{lstlisting}
import Thesis.Prop.BridgeToFoundation

namespace Thesis.Prop

abbrev Theory : Type := LO.Propositional.Theory String

abbrev Provable (T : Theory) (ψ : F) : Prop :=
  T ⊢ ψ

theorem provable_tr_of_tautology (φ : Formula) :
    IsTautology φ → Provable (∅ : Theory) (tr φ) := by
  intro hTaut
  apply LO.Propositional.ClassicalSemantics.completeness!
  intro v hvT
  have : eval v φ := hTaut v
  exact (eval_tr v φ).1 this

end Thesis.Prop
\end{lstlisting}


\section{The Provability Interface: A Natural Deduction System}
\label{sec:nd-interface}

\subsection{Defining the Surface Proof System}
We now define the surface provability relation as a \textbf{Natural Deduction system with contexts and discharge}. This matches the standard pedagogical presentation of logic and supports standard proof techniques like hypothetical reasoning ($\Imp$-Introduction) and classical reductio.

The definition \texttt{ND} is given by an inductive predicate \texttt{ND : Set Formula → Formula → Prop}, representing the judgment $\Gamma \vdash \varphi$.

\begin{lstlisting}
open Set

/-- 
We encode ⊥ inside the {~, ⟶}-only surface language as ~(P ⟶ P).
This is semantically always false, so it behaves as falsity.
-/
def Bot : Formula := ~(.atom "⊥" ⟶ .atom "⊥")

/-- Natural Deduction with contexts (Γ) and discharge. -/
inductive ND : Set Formula → Formula → Prop
| hyp {Γ φ} : φ ∈ Γ → ND Γ φ

| impI {Γ φ ψ} : ND (insert φ Γ) ψ → ND Γ (φ ⟶ ψ)
| impE {Γ φ ψ} : ND Γ (φ ⟶ ψ) → ND Γ φ → ND Γ ψ

| negI {Γ φ} : ND (insert φ Γ) Bot → ND Γ (~φ)
| negE {Γ φ} : ND Γ (~φ) → ND Γ φ → ND Γ Bot

| botE {Γ φ} : ND Γ Bot → ND Γ φ

/-- Classical reductio: from Γ,¬φ ⊢ ⊥ infer Γ ⊢ φ. -/
| classical {Γ φ} : ND (insert (~φ) Γ) Bot → ND Γ φ

/-- The Verified Oracle (The Bridge) -/
| oracle {φ} :
    Provable (∅ : Theory) (tr φ) → ND (∅ : Set Formula) φ
\end{lstlisting}

\subsection{The Architectural Role of the Oracle (Hybrid Proof System)}
The constructor \texttt{oracle} is the defining feature of this architecture. In a traditional ``monolithic'' completeness proof, one would prove that every tautology can be derived using only the standard introduction and elimination rules.

However, formalizing a full completeness proof for Natural Deduction from scratch involves complex syntactic manipulations (e.g., proving that $\Gamma \cup \{\varphi\} \vdash \bot \implies \Gamma \vdash \neg \varphi$ is admissible without using the rule itself). Instead, we use a \textbf{modular interface approach} that results in a \textbf{hybrid proof system}:
\begin{enumerate}
    \item We treat the \emph{Foundation} library as a trusted backend logic engine (providing semantic tableaux).
    \item We treat the surface system \texttt{ND} as a user-facing interface for manual proof construction.
    \item The \texttt{oracle} acts as a \textbf{verified semantic bridge}: it allows the surface system to import complete proofs from the backend, but only at the empty context ($\emptyset$).
\end{enumerate}
This design formally guarantees that any tautology is provable in our system, either by manual derivation or by invoking the backend engine.

\subsection{Completeness Statement}
With the oracle in place, we obtain immediate completeness for the natural deduction system.

\begin{theorem}[ND Completeness via Backend]
For every surface formula $\varphi$,
\[
\mathrm{Taut}(\varphi) \Longrightarrow \emptyset \vdash_{\mathrm{ND}} \varphi.
\]
\end{theorem}
\begin{proof}
    Assume $\mathrm{Taut}(\varphi)$.
    \begin{enumerate}
        \item By the Bridge Lemma (Lemma~\ref{lem:bridge}), we have $\forall v, v \Sat t(\varphi)$.
        \item By the Foundation library's completeness theorem, $\emptyset \Prov t(\varphi)$.
        \item By applying the constructor \texttt{ND.oracle}, we conclude $\mathrm{ND}(\emptyset, \varphi)$.
    \end{enumerate}
\end{proof}

\subsection{The Path to Internalization}
This subsection has been superseded by the detailed discussion in Section~\ref{sec:limitations}.

\section{Limitations and Future Work: Towards Syntactic Completeness}
\label{sec:limitations}

The current system is ``complete by import'': if $\varphi$ is valid, there exists a derivation in \texttt{ND} because the \texttt{oracle} rule allows valid formulas to be accepted as axioms.

A fully internal completeness theorem would remove the reliance on the oracle by proving a \textbf{Syntactic Compilation Theorem}:
\[
(\emptyset \Prov_{\mathrm{Foundation}} t(\varphi)) \implies (\emptyset \vdash_{\mathrm{ND}} \varphi)
\]
This would require implementing a compiler that translates every proof step from Foundation's sequent calculus into a corresponding tree of natural deduction rules. Such a result---establishing that the oracle is \emph{admissible} rather than primitive---is the gold standard of proof-theoretic verification but is significantly more labor-intensive.

The present work establishes the architecture for such a future development: the surface syntax, semantics, and bridge lemma are fully verified, leaving the internal compiler as a modular component for future work.

\section{Main example file (\texttt{Main.lean})}
\label{sec:main}

\subsection{Proof sketch}
This file demonstrates the end-to-end usage:
\begin{enumerate}
\item Define $P$ and $\varphi := (P\Imp P)$.
\item Prove $\mathrm{Taut}(\varphi)$ by simplification.
\item Apply completeness to obtain a derivation in the ND system (via the oracle).
\end{enumerate}

\subsection{Lean code (example)}
\begin{lstlisting}
import Thesis.Prop.NaturalDeduction

namespace Thesis.Prop

def P : Formula := .atom "P"
def taut1 : Formula := P ⟶ P

theorem taut1_is_taut : IsTautology taut1 := by
  intro v
  simp [taut1, eval]

theorem taut1_complete_ND : ND (∅ : Set Formula) taut1 :=
  completeness_ND taut1 taut1_is_taut

end Thesis.Prop
\end{lstlisting}

\section{Conclusion}
\label{sec:conclusion}

The chapter establishes a verified pipeline from semantics to natural deduction:
\[
\mathrm{Taut}(\varphi)
\;\xRightarrow{\text{bridge lemma}}\;
\forall v,\ v\Sat t(\varphi)
\;\xRightarrow{\text{library completeness}}\;
\emptyset\Prov t(\varphi)
\;\xRightarrow{\text{ND oracle}}\;
\mathrm{ND}(\emptyset, \varphi).
\]
The central mathematical work performed locally is the semantics-preservation bridge lemma.
The framework cleanly separates the user-friendly surface system (ND with contexts) from the powerful backend engine (Foundation's semantic tableaux), providing a robust architecture for educational proof assistants.

\begin{thebibliography}{99}

\bibitem{Stansifer2001}
Ryan Stansifer.
\newblock \emph{Completeness of Propositional Logic as a Program}.
\newblock March 2001.

\bibitem{Gentzen1935}
Gerhard Gentzen.
\newblock Untersuchungen "uber das logische Schlie{\ss}en. I.
\newblock \emph{Mathematische Zeitschrift} 35:176--210, 1935.

\bibitem{Prawitz1965}
Dag Prawitz.
\newblock \emph{Natural Deduction: A Proof-Theoretical Study}.
\newblock Almqvist \& Wiksell, 1965. Reprinted by Dover, 2006.

\bibitem{TroelstraSchwichtenberg1996}
A.~S. Troelstra and H.~Schwichtenberg.
\newblock \emph{Basic Proof Theory}.
\newblock Cambridge Tracts in Theoretical Computer Science 43, Cambridge University Press, 1996.

\bibitem{SorensenUrzyczyn2006}
M.~H. S{\o}rensen and P.~Urzyczyn.
\newblock \emph{Lectures on the Curry--Howard Isomorphism}.
\newblock Studies in Logic and the Foundations of Mathematics 149, Elsevier, 2006.

\bibitem{PierceTAPL2002}
Benjamin C.~Pierce.
\newblock \emph{Types and Programming Languages}.
\newblock MIT Press, 2002.

\bibitem{FoundationRepo}
FormalizedFormalLogic.
\newblock \emph{Foundation: A Lean 4 library for formal logic}.
\newblock \url{https://github.com/FormalizedFormalLogic/Foundation}.
\newblock Accessed January 2026.

\end{thebibliography}

\end{document}
